\section*{Scope and Interpretive Limits}

\subsection*{Status of This Work}
ESQET is a \textbf{speculative EFT framework}. It is not presented as a UV-complete theory of quantum gravity. The following limits apply throughout this series:

\begin{itemize}
    \item \textbf{RG equations} are derived within a truncated polynomial approximation and should not be interpreted as scheme-independent RG invariants.
    \item \textbf{Numerical exponents} (e.g., spectral tilts, fixed point values) are illustrative within this truncation and subject to significant uncertainty.
    \item \textbf{No UV completion} is claimed or implied. The framework is effective only up to the highest probed scale ($\mu = 10^{19}$ GeV in our numerical integration).
    \item \textbf{Golden-ratio structure} ($\phi = (1+\sqrt{5})/2$) serves as a \textit{phenomenological organizing motif} motivated by discrete self-similarity, not as an emergent constant or derived universal eigenvalue.
    \item \textbf{Observational estimates} are order-of-magnitude phenomenology intended to illustrate potential experimental signatures, not precise predictions.
\end{itemize}

\subsection*{Geometric Conventions}
The torsion sector follows the Einstein-Cartan formulation with:
\begin{itemize}
    \item Metric compatibility ($\nabla_\mu g_{\nu\rho} = 0$)
    \item Torsion tensor $T^\lambda_{\mu\nu} = \Gamma^\lambda_{[\mu\nu]}$
    \item Parity-odd contributions from the dual torsion tensor $\tilde{T}^{\mu\nu\rho} = \epsilon^{\mu\nu\rho\sigma} T_\sigma$
\end{itemize}
Propagating torsion degrees of freedom are assumed to be suppressed by the scale $\Lambda$.

\subsection*{Golden Ratio as Motif, Not Derivation}
The golden ratio $\phi$ appears as a heuristic organizing principle. Possible origins could include:
\begin{itemize}
    \item Discrete scale invariance with preferred ratio $\phi$
    \item Substitution-system scaling constants (e.g., Tribonacci)
    \item Log-periodic RG spirals with complex critical exponents
\end{itemize}
No derivation uniquely selects $\phi$ over other algebraic irrationals (silver ratio, plastic number, etc.). The framework would remain conceptually similar with any Pisot number as scaling factor.

\subsection*{Beta Function Status}
The RG equation
\[
\mu \frac{d\xi}{d\mu} = -\alpha_1 \xi + \alpha_2 \xi^3 + \alpha_3 \xi C_\alpha
\]
is a \textbf{phenomenological polynomial truncation}, not a derived beta function from explicit loop integrals. A full two-loop calculation in a scheme-independent regularization would be required for quantitative precision.

\subsection*{Spectral Index Caveat}
The predicted tensor tilt $n_T \approx 4.90$ is extraordinarily blue. This value should be interpreted as:
\begin{itemize}
    \item A \textit{transient effective exponent} within a specific truncation
    \item Not an asymptotic inflationary tensor tilt
    \item Likely suppressed by a cutoff or narrow-band mechanism
\end{itemize}
Standard inflationary constraints ($|n_T| \ll 1$) apply to persistent scale-invariant spectra, not necessarily to transient features in truncated EFTs.

\subsection*{Observational Estimate Caveats}
\begin{itemize}
    \item The LISA delay estimate $\Delta t \sim 10^3$ s assumes propagation over cosmological distances without absorption. This should be checked against constraints from GW170817 ($|c_T/c - 1| < 5\times10^{-16}$).
    \item The atom interferometry modulation uses Earth's orbital velocity and rotation; terrestrial laboratory constraints are already stringent ($\xi^2 (T_0/\Lambda)^2 < 10^{-14}$).
    \item All observational forecasts are order-of-magnitude illustrations, not precise experimental predictions.
\end{itemize}

\subsection*{Summary of Interpretive Boundaries}
\begin{tabular}{|p{0.4\textwidth}|p{0.5\textwidth}|}
\hline
\textbf{Claim} & \textbf{Status} \\
\hline
ESQET is a speculative EFT ansatz & ✅ Acknowledged \\
\hline
RG fixed points are truncation-dependent & ✅ Explicitly stated \\
\hline
No UV completion is claimed & ✅ Explicit \\
\hline
Golden ratio is a heuristic motif, not derived & ✅ Explicit \\
\hline
Beta function is a polynomial truncation & ✅ Explicit \\
\hline
Spectral index is transient/effective & ✅ Explicit \\
\hline
Observational estimates are order-of-magnitude & ✅ Explicit \\
\hline
\end{tabular}

\subsection*{Conclusion on Interpretive Status}
This series is presented as \textbf{exploratory theoretical phenomenology}. The mathematical structures (discrete scale invariance, substitution systems, Pisot numbers) are rigorous. Their application to torsion-scalar EFTs and RG flows is speculative. The observational signatures are falsifiable but should be viewed as order-of-magnitude illustrations, not precise predictions.

The value of this work lies in the \textit{organizational arc}:
\begin{center}
EFT ansatz → symmetry breaking → RG truncation → phenomenology → falsifiable tests
\end{center}
not in claims of UV completion or exact golden-ratio derivation.
\end{document}
